\documentclass[a4paper]{article}
\usepackage{amsfonts}
\usepackage{amsmath}
\usepackage{amssymb}
\usepackage{graphicx}     

\begin{document}
  
\noindent \large Auteur: Nathan De Roy \\
\noindent \large Studierichting: Informatica\\
\noindent \large Oefeningenreeks 2D nr. 16.6\\

\medskip

\normalsize

Los op: $f(x) = \dfrac{1 + \cos(2x)}{\sin(x)}$ voor: $x \to \dfrac{\pi}{2}; x \to 0; x \to \pi$\\

Voor $x \to \dfrac{\pi}{2}: \lim\limits_{x \to \tfrac{\pi}{2}}\left(\dfrac{1 + \cos(2x)}{\sin(x)}\right) = \dfrac{1 + \cos(2\dfrac{\pi}{2})}{\sin(\dfrac{\pi}{2})} = \dfrac{1 + -1}{1} = 0$\\


Voor $x \to 0: \lim\limits_{x \to 0}\left(\dfrac{1 + \cos(2x)}{\sin(x)}\right) = \dfrac{1 + \cos(0)}{\sin(0)} = \dfrac{1 + 1}{0} = \infty$\\

Bekijk linker- en rechterlimiet:\\

$\indent \lim\limits_{x \to 0+}\left(\dfrac{1 + \cos(2x)}{\sin(x)}\right) = \dfrac{1 + \cos(0+)}{\sin(0+)} = \dfrac{1 + 1+}{0+} = +\infty$\\

$\indent \lim\limits_{x \to 0-}\left(\dfrac{1 + \cos(2x)}{\sin(x)}\right) = \dfrac{1 + \cos(0-)}{\sin(0-)} = \dfrac{1 + 1+}{0-} = -\infty$\\

Voor $x \to \pi: \lim\limits_{x \to \pi}\left(\dfrac{1 + \cos(2x)}{\sin(x)}\right) = \dfrac{1 + \cos(2\pi)}{\sin(\pi)} = \dfrac{1 + 1}{0} = \infty$\\

Bekijk linker- en rechterlimiet:\\

$\indent \lim\limits_{x \to \pi+}\left(\dfrac{1 + \cos(2x)}{\sin(x)}\right) = \dfrac{1 + \cos(2\pi+)}{\sin(\pi+)} = \dfrac{1 + 1+}{0-} = -\infty$\\

$\indent \lim\limits_{x \to \pi-}\left(\dfrac{1 + \cos(2x)}{\sin(x)}\right) = \dfrac{1 + \cos(2\pi-)}{\sin(\pi-)} = \dfrac{1 + 1+}{0+} = +\infty$\\

\begin{thebibliography}{999}
%\bibitem{a} Referentie
\end{thebibliography}
\end{document} 